% Six wires from the Arduino header to the ADXL345 breakout, plus the
% interrupt line. 3.3 V only.
\begin{circuitikz}[font=\sffamily\scriptsize]

\node[draw=linecol, fill=hwcol, minimum width=32mm, minimum height=62mm,
      align=center, rounded corners=2pt] (nuc) at (0,-0.2) {};
\node[font=\sffamily\bfseries\scriptsize, align=center] at (0,2.5)
  {NUCLEO-H7A3ZI-Q\\{\tiny Arduino header}};

\node[draw=linecol, fill=extcol, minimum width=30mm, minimum height=54mm,
      align=center, rounded corners=2pt] (brk) at (8.0,-0.4) {};
\node[font=\sffamily\bfseries\scriptsize, align=center] at (8.0,2.1)
  {SEN0032\\{\tiny ADXL345 breakout}};

\foreach \y/\l/\r/\sig in {
   1.4/{3V3}/{VCC}/{supply 3.3 V},
   0.7/{D13}/{SCK}/{serial clock},
   0.0/{D11}/{SDI}/{master out slave in},
  -0.7/{D12}/{SDO}/{master in slave out},
  -1.4/{D10}/{CS}/{chip select},
  -2.1/{D2}/{INT1}/{watermark interrupt},
  -2.8/{GND}/{GND}/{common return}}{
  \node[font=\tiny, anchor=west] at (-1.5,\y) {\l};
  \node[font=\tiny, anchor=west] at (6.65,\y) {\r};
  \draw[cable] (1.6,\y) -- (6.5,\y);
  \node[font=\tiny, fill=white, inner sep=1.2pt] at (4.05,\y+0.21) {\sig};
}

% the common return, dropped from the lowest wire so it crosses nothing
\draw[cable] (4.05,-2.8) -- (4.05,-3.7);
\node[ground] at (4.05,-3.7) {};

\node[umlnote, text width=54mm, anchor=north west] at (-1.9,-4.5)
  {Nothing here sees more than 3.3 V and the 5V pin is not used. Nothing on
   this bench can be soldered, so if the breakout's bus selection turns out to
   need a bridge moved, the chapter uses the other bus instead.};
\node[umlnote, text width=46mm, anchor=north west] at (5.2,-4.5)
  {Every header pin above is the Nucleo-144 convention until it has been read
   in the board manual. That is confirm item 12.};
\end{circuitikz}
